\section{The arithmetic substrate}\label{sec:arithmetic}

\subsection{Integer, rational, and complex-rational codes}

An integer is represented by a pair \((u,v)\in\N^2\), denoting \(u-v\).
A rational is a pair \((m,d)\), where \(m\) is an integer code and \(d+1\)
is the positive denominator.  Equality and order are tested by cross
multiplication.  Normalization by the bounded Euclidean algorithm is useful
for canonical input codes but is not mathematically necessary.  A complex
rational is a pair of rational codes.  We use
\[
  \norm{x+iy}=|x|+|y|
\]
for error estimates.  It satisfies
\(\norm{uv}\le\norm u\norm v\), and every comparison is rational.
The Euclidean norm is used only in familiar displayed formulas; replacing it
by the one-norm changes constants but not the constructions.

Division is totalized.  Define
\[
 \operatorname{qdiv}(x,y)=
 \begin{cases}x/y,&y\ne0,\\0,&y=0.
 \end{cases}
\]
Every theorem that uses the first branch carries a proof of \(y\ne0\).  This
allows all displayed functions to be primitive recursive without turning a
partial operation into an implicit logical principle.

\subsection{Finite sequences and bounded folds}

Pairing and projections code finite tuples.  Bounded sums, products, maxima,
and Boolean folds are obtained by primitive recursion:
\begin{align*}
 \operatorname{sum}(0,f)&=0,&
 \operatorname{sum}(n+1,f)&=\operatorname{sum}(n,f)+f(n),\\
 \operatorname{prod}(0,f)&=1,&
 \operatorname{prod}(n+1,f)&=\operatorname{prod}(n,f)f(n).
\end{align*}
A bounded universal assertion is represented by a Boolean fold
\[
  \operatorname{all}(N,P)=\prod_{n<N}\chi_P(n).
\]
It is equal to \(1\) precisely when every decidable test \(P(n)\) below \(N\)
holds.  This is the standard way to keep induction invariants literally
quantifier-free.  For example, instead of carrying
\(\forall k<n\,D_k\ne0\), the object-theory proof carries
\(\operatorname{all}(n,k\mapsto[D_k\ne0])=1\).

\subsection{Names of effective numbers}

\begin{definition}[Dyadic Cauchy name]\label{def:name}
A complex number name is a primitive recursive function
\(X:\N\to\Q(i)\) satisfying
\[
  \norm{X(p+1)-X(p)}\le2^{-p}
  \qquad(p\in\N).
\]
A stronger and often more convenient contract is
\[
  \norm{X(q)-X(p)}\le2^{-p+1}
  \qquad(q\ge p).
\]
Two names are extensionally equal when a primitive recursive shift \(s\)
satisfies
\[
  \norm{X(p+s)-Y(p+s)}\le 2^{-p}
  \quad\text{for every }p.
\]
\end{definition}

The shift absorbs a finite number of error sources.  Given a rational
constant \(C\ge0\), define by bounded search
\[
 \sigma(C)=\min\{s\le \lceil C\rceil+1:2^s\ge\max(1,C)\}.
\]
Existence follows from \(2^s\ge s+1\), proved by quantifier-free induction.
Then \(C2^{-p-\sigma(C)}\le2^{-p}\).  Every precision allocation below is a
composition of such shifts, bounded searches, and elementary arithmetic, and
is therefore primitive recursive.

\begin{definition}[Named function on a rational domain]\label{def:named-function}
A named function on a coded rational domain \(D\) consists of a primitive
recursive evaluator
\[
  F(x,p)\in\Q(i).
\]
A uniform value bound \(M\), and a modulus \(\omega\) such that
\[
  x,y\in D,\quad \norm{x-y}\le2^{-\omega(p)}
   \Longrightarrow \norm{F(x,p+s)-F(y,p+s)}\le2^{-p}
\]
for a displayed shift \(s\).  Derivative, integral, and branch contracts add
further finite inequalities.  The evaluator is fixed; the theory does not
quantify over all evaluators.
\end{definition}

\subsection{Apartness and branch data}

Constructive division requires more than the negation of equality.  To name
\(1/X\), supply a positive rational margin \(\delta\) and a precision
\(p_0\) such that
\[
  \norm{X(p_0)}\ge\delta,
  \qquad
  \norm{X(p)-X(p_0)}\le\delta/2\quad(p\ge p_0).
\]
Then \(\norm{X}\ge\delta/2\), and rational interval arithmetic gives a name
for the reciprocal.  We call this an \emph{apartness certificate}.

A branch of \(z^\alpha\) on a rational disk consists of a named logarithm
\(L(z)\), a disk center and radius, a proof that the disk avoids the branch
cut and zero by a rational margin, and the local identities
\(L'(z)=1/z\), \(L(z_0)=\ell_0\).  Then
\(z^\alpha=\exp(\alpha L(z))\).  The exponential is represented by its
Taylor series with a factorial-tail modulus.  No global choice of argument is
assumed.

\subsection{Formal power series as finite observations}

A full infinite sequence is not an object over which first-order iPRA
quantifies.  A fixed primitive recursive coefficient function
\(A:\N\to\Q(i)\), however, is legitimate.  A formal-series identity
\[
  \sum_{n\ge0}A_nz^n=\sum_{n\ge0}B_nz^n
\]
means the theorem schema \(A_n=B_n\) for every \(n\).  Product and
composition identities are expressed through finite Cauchy sums.  For
example, the coefficient of a product is
\[
  (A*B)_n=\sum_{k=0}^{n}A_kB_{n-k}.
\]
Every coefficient uses only a bounded sum.

\begin{definition}[Truncated series algebra]\label{def:trunc-series}
For \(N\in\N\), let
\(R_N=R[z]/(z^N)\).  Represent an element by the vector
\((a_0,\ldots,a_{N-1})\).  Addition, Cauchy multiplication, formal
differentiation, and substitution by a series with zero constant term are
primitive recursive operations on finite vectors.  Equality in \(R_N\) is a
finite conjunction.
\end{definition}

A uniform formal identity is proved by an induction whose invariant is the
correctness of all coefficients below the current bound.  Alternatively, for
a fixed target coefficient \(n\), one works entirely in \(R_{n+1}\).  This
``finite observation'' interpretation is especially useful for the Kummer
forms and algebraic substitutions.

\subsection{Isolated algebraic names}

Radicals are handled without postulating a total root operation on arbitrary
reals.

\begin{definition}[Isolated algebraic number]\label{def:alg-name}
An isolated real algebraic name consists of an integer polynomial \(P\),
rational endpoints \(u<v\), a sign-change or monotonicity certificate proving
that \(P\) has exactly one zero in \([u,v]\), and a primitive recursive
bisection procedure.  A complex algebraic name uses two real intervals plus a
finite uniqueness certificate, for example a rectangle on which a resultant
or Newton test is separating.
\end{definition}

Bisection has a primitive recursive convergence rate because the interval
width after \(m\) steps is \((v-u)2^{-m}\).  Square and cube roots occurring
in the source pages are interpreted by positive isolated roots.  Polynomial
identities between such names are proved by interval refinement together with
an exact polynomial relation, not by an unguarded appeal to classical root
choice.

\subsection{Integration by finite quadrature}

For a named uniformly continuous function \(f:[0,1]\to\C\), define the
left Riemann sum
\[
  R_N(f)=\frac1N\sum_{j<N}f(j/N).
\]
If \(\norm{f(x)-f(y)}\le\varepsilon\) whenever
\(|x-y|\le1/N\), then
\[
  \norm{R_N(f)-R_{2N}(f)}\le\varepsilon.
\]
The proof pairs the two fine cells in every coarse cell and uses the supplied
modulus.  Thus a primitive recursive mesh schedule produces a Cauchy name for
the integral.  Endpoint singularities are treated by cutting off
\([0,\eta]\) and \([1-\eta,1]\), bounding those tails by explicit powers,
and applying ordinary quadrature on the central interval.  This is the
integration engine used in \cref{sec:integrals-sums}.

\begin{lemma}[Finite substitution principle]\label{lem:finite-substitution}
Suppose \(f,g\) are named functions on rational intervals, \(g\) is
monotone with named derivative and explicit endpoint values, and the
fundamental theorem estimate is supplied with a modulus.  Then the usual
change-of-variables formula is an iPRA theorem schema.
\end{lemma}

\begin{proof}
Apply the identity to the polygonal/Riemann sums.  On each cell, replace the
increment \(g(x_{j+1})-g(x_j)\) by
\(g'(x_j)(x_{j+1}-x_j)\).  The derivative-remainder contract bounds the
sum of errors by a primitive recursive multiple of the mesh.  The two finite
sums telescope at the endpoints.  Choosing the mesh by bounded search makes
the total error at most \(2^{-p}\).  Every invariant is an equality or
inequality of finite rational sums.
\end{proof}
